\section{Robustness and interpretation}\label{sec:short_robustness}

The findings are similar under alternative assumptions about historical returns. A standardized Student-\(t\) model leaves the PI--PM difference small relative to historical-volatility pricing errors. Reconstructing first-nearby CL returns contract by contract gives a similar result: on the 2,381 out-of-sample observations, historical PI has MAE 2.6468 and RMSE 3.4451, compared with 2.6187 and 3.4090 using the continuous/front-month return proxy. Further historical-volatility benchmarks appear in Online Supplement Section S1.

The small observed PI--PM difference is larger than the numerical approximation error. Monte Carlo with many draws, independently scrambled Sobol simulation, and Curran's approximation agree closely on the price difference. Simulation-based calibration provides a separate check on the posterior implementation. The supplement also examines pricing at the posterior root-mean-square volatility (Section S2), alternative ways to map LO implied volatility to one APO volatility, and the sensitivity to the number of Cox--Ross--Rubinstein tree steps (Section S4). These numerical checks do not establish an exact arbitrage-free representation of the volatility surface.

The market comparisons have important limitations. In the
original sample of 2,381 option-date observations, 2,319
have zero recorded trading volume.
Earlier option-implied volatility still yields lower errors
than the tested historical estimates when we restrict the
sample to positive-volume or open-interest observations,
but these subsamples are small, especially for some expiries
(Online Supplement Section S3).
An exchange settlement proxy should not be interpreted
as an executable transaction price. Moreover, the main pricing comparison contains only 15 target dates and the external validation only 10. Resampling whole valuation dates accounts for dependence among contracts on the same date, but may not fully account for dependence between nearby dates.

Finally, the one-factor constant-volatility model excludes richer futures-curve dynamics, stochastic volatility, jumps, and convenience yields. The implied volatilities are values inferred under the chosen pricing model. They should not be interpreted as uniquely identified risk-neutral diffusion volatilities or as direct estimates of a variance risk premium. The reported pricing-error differences apply to these volatility inputs and this pricing model; they do not separate physical and risk-neutral sources of volatility structurally.

The Databento analysis covers more dates but has different
measurement limitations. Forward averages and effective discount
factors are inferred from other \emph{same-date} option strikes;
the historical volatility proxy is estimated from earlier
parity-implied forward-average changes, not first-nearby CL
realized returns; and the lognormal terminal-average approximation
replaces the path of known and remaining daily fixings
with a model for the terminal average. We cannot independently confirm that the original Barchart
end-of-day \texttt{Latest} fields are final exchange
settlements or executable trade prices.
CME Clearing switched AAO end-of-day settlement valuations to
its Curran model effective May 18, 2020
\citep{cmeclearing2020curran}. Agreement with end-of-day settlement prices may therefore
reflect the stability or reproduction of exchange model
valuations rather than new information about tradable prices.

A very small price correction need not have economic value.
For instance, the older October-expiry mean absolute PI--PM
difference is about \$0.000745/bbl, well below the standard
\$0.01/bbl quotation increment. In the longer conditional comparison, Bayesian modeling of
implied volatility and simply carrying forward the last
observation have very similar errors. Resampling by month
does not establish a consistent advantage for either method. The lower-than-intended coverage of pricing-error intervals
in early 2026 also cautions against assuming that an uncertainty
model fitted to past errors remains well calibrated over time.

The Barchart/LO and Databento analyses differ in the
observed prices, data sources, pricing models, available
expiries, and information used for valuation. They
are \emph{separate analyses}, not pooled observations for one
common estimator. Neither tests executable hedging profits or measures a
variance risk premium. We still need to verify when the vendors originally published
their records, whether those records were revised, whether
Barchart data may be redistributed, and whether the original
market prices and full input histories can be independently checked.
The arithmetic behind the 2,381-observation fitted-smile
comparison and the 2,366-observation same-contract comparison
has been independently checked, including all 4,762 stored
implied-volatility estimates and the recorded ordering of dates. These checks do not establish when the source records first
became available.

